\chapter{绪论}\label{chap:introduction}

\section{研究背景与意义}

宝可梦对战（Pokémon Battle）是宝可梦系列作品的核心玩法之一。自1996年《宝可梦 红/绿》发售以来，宝可梦对战系统经历了数十年的发展，已形成包含属性克制、特性触发、道具效果、场地状态等数十种要素的复杂博弈体系。在官方对战规则（VGC, Video Game Championships）和民间对战平台（Pokémon Showdown等）的推动下，宝可梦对战已发展成为一项具有高度竞技性的策略游戏。

截至2026年，宝可梦全国图鉴编号已突破1500种，每种宝可梦具有独特的种族值分布、属性组合、特性池和技能池。在6v6单打对战、4v4双打对战等不同赛制下，如何从海量的可能性中筛选出最优的队伍配置和战术策略，是一个具有重要研究价值的组合优化与博弈论交叉问题。

% 引用格式说明：
% \citep{} = 括号引用，如"...已有研究对此进行了探讨(张三, 2020)"
% \citet{} = 正文引用，如"张三(2020)的研究表明..."
% 多处引用用英文逗号分隔：\citep{ref1, ref2, ref3}
近年来，国内外学者开始将计算博弈论、蒙特卡洛树搜索、强化学习等方法引入宝可梦对战研究\citep{lamport1986document,walls2013drought}。然而，现有研究多局限于单一机制的分析，缺乏对属性克制、技能搭配、努力值分配等多维度的综合优化框架。本文旨在构建一个面向宝可梦PVP对战的策略优化模型，为训练家提供科学的队伍配置决策支持。

\begin{figure}[!htbp]
    \centering
    \includegraphics[width=0.45\textwidth]{pokemon/pikachu}
    \bicaption{宝可梦对战已发展为高度复杂的策略博弈体系}{Pokémon battles have evolved into highly complex strategic systems}
    \label{fig:pikachu}
\end{figure}

\section{国内外研究现状}

\subsection{属性克制与队伍覆盖优化}

属性克制是宝可梦对战的基础机制。现有18种属性（包括2013年引入的妖精属性）构成了一个复杂的攻防克制网络。\citet{chen1980zhongguo} 对属性覆盖面问题进行了初步探讨，提出了基于集合覆盖理论的队伍属性配置方案。然而，该研究未考虑宝可梦的双属性组合和特性修正对属性克制的实际影响。

近年来，\citet{betts2005aging} 将图论方法引入属性覆盖分析，将属性克制关系建模为加权有向图，通过求解最小支配集来优化队伍的攻防覆盖。但该方法计算复杂度较高，难以应用于实时对战场景。

\subsection{速度线与努力值分配}

速度线（Speed Tier）是决定先手权的重要因素。\citet{chu2004tushu} 通过统计分析VGC 2022赛季的使用率数据，提炼了关键速度阈值。随后，\citet{walls2013drought} 将努力值分配问题建模为多目标优化问题，利用遗传算法求解，在特定赛制下取得了优于人工分配的结果。然而，该方法存在收敛速度慢、易陷入局部最优等不足。

\subsection{强化学习在宝可梦对战中的应用}

深度强化学习在游戏AI领域取得了显著进展。\citet{hls2012jinji} 使用DQN算法训练宝可梦对战AI，在有限精灵池内取得了超越人类玩家的胜率。但是，这些AI的决策过程缺乏可解释性，难以直接转化为训练家可以理解的战术指导。此外，现有AI模型通常只考虑单打对战场景，对双打对战的协同配合关注不足。

\subsection{现有研究的不足之处}

综合以上分析，现有研究存在以下不足：(1) 属性覆盖优化未充分考虑双属性组合和特性修正的影响；(2) 速度线分析多依赖统计方法，缺乏理论建模；(3) 各优化目标之间缺乏协同；(4) AI决策缺乏可解释性。

\section{本文主要研究内容}

本文围绕宝可梦PVP对战策略优化这一核心问题，从属性覆盖、速度优化和队伍协同三个维度展开研究，主要贡献如下：

\begin{enumerate}
    \item 提出一种融合双属性组合与特性修正的属性覆盖面优化算法。该算法在传统属性克制矩阵的基础上，引入双属性组合效应因子和特性修正项，能够更准确地评估队伍对主流环境宝可梦的克制覆盖率。
    \item 构建基于贝叶斯优化的速度线阈值预测模型。通过对VGC赛事数据的分析，实现对特定赛制下关键速度门槛的高精度估计，为训练家的努力值分配提供定量依据。
    \item 设计面向队伍整体协同的遗传算法框架。该框架将属性覆盖、速度配合和攻防平衡作为多个优化目标，通过非支配排序遗传算法（NSGA-II）求解帕累托最优解集。
    \item 开发宝可梦配招推荐原型系统，通过机器学习模型为给定宝可梦推荐最优的技能组合，并验证了所提方法的有效性。
\end{enumerate}

\section{论文组织结构}

本文共分为六章。第一章为绪论，介绍研究背景、国内外研究现状和主要研究内容。第二章介绍宝可梦对战的基础理论，包括属性克制、种族值、努力值、技能系统等。第三章详细阐述属性覆盖面优化算法及其实验验证。第四章阐述速度线预测模型与队伍协同优化框架的构建。第五章通过大规模对战模拟实验评估所提方法的有效性。第六章总结全文并展望未来研究方向。
